\section*{第 \ref{ch:b3:molecular-evolution} 章——分子演化与系统发生基因组学}

\begin{solution}{exo:b3:molecular-evolution:1}
$k = \mu$：每个位点、每代的中性\omterm{thm:b3:molecular-evolution:neutral}{替换速率}等于
突变速率。每代出现 $2N\mu$ 个新的中性突变，
每个突变的\omterm{thm:b3:molecular-evolution:neutral}{固定概率}为 $1/2N$；种群越大，突变越多，
但每个突变的固定机会按比例减小，
两种效应恰好抵消。
\end{solution}

\begin{solution}{exo:b3:molecular-evolution:2}
\omterm{def:b3:genomics:comparative}{直系同源基因}：两个物种中的基因，二者源自共同祖先中的同一个基因，
例如人类 $\beta$-珠蛋白基因与小鼠 $\beta$-珠蛋白基因。
\omterm{def:b3:genomics:comparative}{旁系同源基因}：同一基因组中由\omterm{def:b3:molecular-evolution:duplication}{基因复制}产生的基因，
例如人类的 $\alpha$-与 $\beta$-珠蛋白基因，或 $\beta$ 与 $\gamma$ 基因。\omterm{def:b3:molecular-evolution:duplication}{假基因}：
已退化、不再编码蛋白质的拷贝，即 $\psi\beta$ 基因，
位于人类 $\beta$-珠蛋白基因簇中，其中仍含有终止密码子。
\end{solution}

\begin{solution}{exo:b3:molecular-evolution:3}
$\omega$ 比较改变氨基酸的\omterm{thm:b3:molecular-evolution:neutral}{替换速率}与
沉默\omterm{thm:b3:molecular-evolution:neutral}{替换速率}之比，后者即中性基线。$0.02$：
强\omterm{def:b3:molecular-evolution:dnds}{纯化选择}，约 \qty{98}{\%} 的氨基酸改变被清除——
肌动蛋白、\omterm{def:b3:chromatin-epigenetics:nucleosome}{组蛋白} H3。$1.0$：无约束，如\omterm{def:b3:molecular-evolution:duplication}{假基因}。$2.5$：
\omterm{def:b3:molecular-evolution:dnds}{正选择}，氨基酸改变受到青睐——
MHC 的肽结合槽、HIV 的\omterm{def:b3:virology:virus}{包膜}、反刍动物的溶菌酶。
\end{solution}

\begin{solution}{exo:b3:molecular-evolution:4}
比较哺乳动物间的蛋白质，每个位点约每十亿年发生一次替换；
对于具有十亿个编码位点的基因组，以及几年一代的世代时间，
这大约相当于每代固定一次替换，或每两年固定一次。
霍尔丹的选择代价允许
大约每三百代固定一次受选择的替换。
两者相差一百倍或更多，
因此几乎所有固定的变化都必须是中性的。
\end{solution}

\begin{solution}{exo:b3:molecular-evolution:5}
中性情形：$1/2N = 10^{-4}$。$s = 0.001$，$4Ns = 20$：$(1 - \mathrm{e}^{-0.002})/
(1 - \mathrm{e}^{-20}) = 0.002$。$s = 0.01$：$0.0198$。$s = -0.001$：$(1 -
\mathrm{e}^{0.002})/(1 - \mathrm{e}^{20}) = 0.002\,\mathrm{e}^{-20} =
4\times 10^{-12}$。这些情形都不是有效中性：要使其为有效中性，需要 $|4Ns|
\lesssim 1$，即 $|s| \lesssim 5\times 10^{-5}$；
选择系数仅为千分之一，
在五千个个体的种群中便已具有决定性作用。
\end{solution}

\begin{solution}{exo:b3:molecular-evolution:6}
未校正：$T = 0.12/(2\times 2\times 10^{-9}) = \qty{30}{Myr}$。
Jukes--康托尔校正：$d = -\tfrac{3}{4}\ln(1 - 0.16) = 0.131$，$T =
\qty{33}{Myr}$。校正量达到两倍时，$-\tfrac{3}{4}
\ln(1 - 4p/3) = 2p$，即 $p \approx 0.6$ 时，
原始差异已丢失大部分信息。
\end{solution}

\begin{solution}{exo:b3:molecular-evolution:7}
$p_{N} = 7/700 = 0.010$，$p_{S} = 10/200 = 0.050$，$\omega = 0.2$：
\omterm{def:b3:molecular-evolution:dnds}{纯化选择}清除了五分之四的氨基酸改变。若取
$35$ 与 $10$：$p_{N} = 0.05$，$\omega = 1$：未检测到约束——
可能是\omterm{def:b3:molecular-evolution:duplication}{假基因}，或是某些位点受\omterm{def:b3:molecular-evolution:dnds}{正选择}、另一些位点受\omterm{def:b3:molecular-evolution:dnds}{纯化选择}的基因，
两种选择相互平衡。
\end{solution}

\begin{solution}{exo:b3:molecular-evolution:8}
$k = d/2T = 0.012/(1.3\times 10^{7}) = 9.2\times 10^{-10}$，每个位点每年；
$\mu = 25k = 2.3\times 10^{-8}$，每个位点每代；$6\times
10^{9}\times 2.3\times 10^{-8} \approx 140$，即每个二倍体基因组每代的新突变数。
直接对亲代与子代测序，只找到约七十个。分歧时间的估计
被祖先种群中已有的多态性抬高，
这一多态性使表观分化时间增加了
几十万代。
\end{solution}

\begin{solution}{exo:b3:molecular-evolution:9}
$\mathrm{e}^{-T/2N} = \mathrm{e}^{-1} = 0.37$；不一致比例为 $\tfrac{2}{3}
\times 0.37 = 0.25$。若要达到 \qty{10}{\%}：$\mathrm{e}^{-T/2N} = 0.15$，$T =
2N\ln(1/0.15) = \num{190000}$ 代。
\end{solution}

\begin{solution}{exo:b3:molecular-evolution:10}
当 $s \to 0$ 时：$1 - \mathrm{e}^{-2s} \approx 2s$，$1 - \mathrm{e}^{-4Ns}
\approx 4Ns$，比值为 $1/2N$。对于 $4Ns \gg 1$，分母为 $1$，
分子 $\approx 2s$。对于 $s < 0$ 且 $4N|s| \gg 1$：分子 $1 -
\mathrm{e}^{2|s|} \approx -2|s|$，分母 $1 - \mathrm{e}^{4N|s|} \approx
-\mathrm{e}^{4N|s|}$，因此 $P \approx 2|s|\,\mathrm{e}^{-4N|s|}$。
在 $N = 10^{4}$ 时，比中性情形低一百倍：$2|s|\,\mathrm{e}^{-4\times 10^{4}|s|} =
5\times 10^{-7}$；$|s| = 1.6\times 10^{-4}$ 时，$3.2\times 10^{-4}\times
\mathrm{e}^{-6.4} = 5.3\times 10^{-7}$。因此 $|s| \approx 1.6\times 10^{-4}$，
$4N|s| \approx 6.5$：万分之一的劣势，
在一万个个体中就足以
使\omterm{thm:b3:molecular-evolution:neutral}{固定概率}比随机情形低一百倍。
\end{solution}

\begin{solution}{exo:b3:molecular-evolution:11}
设 $a$ 为分歧点到 $O$ 的距离，两条路径上均如此，则 $a + m =
0.30$，$a + h = 0.24$，$m + h = 0.20$：$m - h = 0.06$，故 $m = 0.13$，$h
= 0.07$，比值为 $1.9$。若采用严格的每代\omterm{prop:b3:molecular-evolution:clock}{分子钟}，
比值应为世代数之比，即 $25/0.5 = 50$。观测到的 $1.9$ 意味着，
小鼠每年积累的变化仅为人类的两倍，
尽管其世代数多出五十倍：每代突变速率必定
在人类中高约 $25$ 倍（即每代生殖系细胞分裂更多），
从而使两谱系每年的差异仅相差
两倍。
\end{solution}

\begin{solution}{exo:b3:molecular-evolution:12}
产生终止密码子的替换以 $1000\times 2\times 10^{-9}\times
0.04 = 8\times 10^{-8}$ 的速率每年出现：首次出现的期望时间约为
\qty{12}{Myr}（移码突变大约使这一时间减半）。在这段时间内，
有利突变（占全部突变的 $10^{3}$ 至 $10^{5}$ 分之一）
远比终止密码子出现的可能性小，即便出现，
其\omterm{thm:b3:molecular-evolution:neutral}{固定概率}也仅为 $2s$。因此，该拷贝通常在
选择找到其用途之前便已失活——脊椎动物中重复基因观测到的半衰期为
数百万年。
\end{solution}

\begin{solution}{pb:b3:molecular-evolution:1}
\textbf{1.} $k = 0.012/(2\times 6.5\times 10^{6}) = 9.2\times 10^{-10}$ 每
位点每年；$\mu = 25k = 2.3\times 10^{-8}$ 每位点每代。
\textbf{2.} $6\times 10^{9}\times 2.3\times 10^{-8} \approx 140$ 个新
突变每个二倍体基因组；其中 \qty{1.5}{\%}，约 $2$ 个，位于
编码序列中。
\textbf{3.} 每个单倍体基因组每代固定的替换数约为 $3\times 10^{9}\times 2.3\times 10^{-8}
\approx 70$ 个，而霍尔丹的数值为
$1/300$：多出两万倍，远超选择所能驱动的量，故
绝大多数为中性。
\textbf{4.} $2N\mu = 10^{5}\times 2.3\times 10^{-8} = 2.3\times 10^{-3}$
每位点每代在种群中产生的新突变；其中 $1/2N =
10^{-5}$ 会被固定。
\textbf{5.} $4N = \num{200000}$ 代，即 \qty{5}{Myr}——几乎与
分化以来的时间相当。
\textbf{6.} 分歧度统计的是自共同祖先以来，沿两条
独立谱系产生的突变；每个突变都在其所在谱系内固定，
长期速率为 $\mu$，与每个突变所需的时间无关，
因为突变持续稳定地产生，滞后只会延迟、
不改变其流动。（唯一需要修正的是分化时存在的多态性：
它使两个等位基因的共同祖先早于
物种分化。）
\textbf{7.} 中性 $10^{-5}$；$s = 0.001$：$0.002$；$s = 0.01$：$0.0198$；
$s = 0.1$: $1 - \mathrm{e}^{-0.2} = 0.18$.
\textbf{8.} $s = 1/4N = 5\times 10^{-6}$：低于此值时，漂变压倒
选择，突变表现为中性；高于此值时，选择
决定其固定的概率。
\textbf{9.} $s = -0.001$：$P = 0.002\,\mathrm{e}^{-200}$，就一切实际目的而言为零
——仅为中性的 $10^{-85}$。$s = -10^{-5}$，$4N|s| = 2$：$P =
2\times 10^{-5}/(\mathrm{e}^{2} - 1) = 3.1\times 10^{-6}$，为中性的 \qty{31}{\%}
：轻度有害突变确实会被固定。
\textbf{10.} 每位点每代的有利替换：$2N\times
10^{-5}\mu\times 2s = 10^{5}\times 10^{-5}\times 0.02\,\mu = 0.02\mu$;
中性：$\approx\mu$。适应性比例为 $0.02/1.02 \approx \qty{2}{\%}$。
尽管每个有利突变被固定的可能性是
中性突变的两千倍，但它们如此罕见，以至于只有五十分之一的替换
是适应性的：选择确实在起作用，却几乎不显现。
\textbf{11.} $p_{N} = 46/920 = 0.050$，$p_{S} = 84/280 = 0.30$。校正后：
$d_{N} = -\tfrac{3}{4}\ln(1 - 0.0667) = 0.052$, $d_{S} = -\tfrac{3}{4}
\ln(1 - 0.40) = 0.38$. $\omega = 0.14$.
\textbf{12.} \omterm{def:b3:molecular-evolution:dnds}{纯化选择}清除了约 \qty{86}{\%} 的
氨基酸变化——一个典型的基因。$k = d_{S}/2T = 0.38/(1.8\times
10^{8}) = 2.1\times 10^{-9}$ 每位点每年：是人–黑猩猩
数值的两倍，这符合预期，因为该比较包含了替换较快的
啮齿类谱系；两者处于同一数量级。
\textbf{13.} $d = -\ln(1 - 0.55) = 0.80$ 每位点；$T = 0.80/(2\times
10^{-9}) = \qty{400}{Myr}$.
\textbf{14.} 两条不同的肽链组成 $\alpha_{2}\beta_{2}$
四聚体，其协同的氧结合、\omterm{prop:b3:structural-biology:mechanism}{玻尔效应}以及
由 \omterm{prop:b3:structural-biology:mechanism}{2,3-二磷酸甘油酸}实现的别构调控，都依赖于
不同亚基之间的界面；而 $\beta$ 基因簇随后可分化为
胚胎型、胎儿型和成体型链，它们具有不同的亲和力。
\textbf{15.} $2kT = 2\times 8\times 10^{-9}\times 4\times 10^{7} = 0.64$
每位点替换；考虑到二十种氨基酸之间的饱和，
原始差异约为 $0.95(1 - \mathrm{e}^{-0.64/0.95}) \approx 0.47$。
在 \qty{500}{Myr} 时，$d = 8$：每个位点都已多次变化，
原始差异达到其上限 \qty{95}{\%}，不再携带
任何关于时间的信息。
\textbf{16.} $10^{-11}\times 10^{9}\times 100\times 2 = 2$ 个替换，
两条谱系在十亿年间的 100 个位点上。对于
\qty{10}{Myr} 的分化，期望值为 $0.02$：几乎总是零差异，
无法分辨。
\textbf{17.} $1200\times 2\times 10^{-9}\times 0.04 = 9.6\times 10^{-8}$
每年：首个终止密码子的期望出现时间约为 \qty{10}{Myr}。
\textbf{18.} 全基因组加倍会同时使每个基因加倍，
因此相互作用蛋白的化学计量得以保持，不存在剂量失衡
而对拷贝进行选择；单个串联重复只改变
一个基因的剂量，往往有害，或者立即变得冗余，
从而可以自由衰变。
\textbf{19.} $\mathrm{e}^{-\num{80000}/\num{100000}} = \mathrm{e}^{-0.8} = 0.45$.
\textbf{20.} $\tfrac{2}{3}\times 0.45 = 0.30$：与
观测到的 \qty{30}{\%} 相符。
\textbf{21.} 各占一半，即 \qty{15}{\%} 的所有树将人与
大猩猩归为一组，\qty{15}{\%} 将黑猩猩与大猩猩归为一组。若其中一种明显偏多，
则表明两物种分化之后存在基因流，
而漂变无法产生这一结果。
\textbf{22.} $N = \num{10000}$：$\mathrm{e}^{-4} = 0.018$，不一致率
\qty{1.2}{\%}。因此，观测到的 \qty{30}{\%} 要求祖先种群规模
约为五万——是当今人类有效
种群大小的数倍。
\textbf{23.} 串联拼接法假定每个基因都具有\omterm{prop:b3:molecular-evolution:ils}{物种树}；
在\omterm{prop:b3:molecular-evolution:ils}{不完全谱系分选}的情况下，各基因具有多种树，而对于较短的
内部分支，最常见的\omterm{prop:b3:molecular-evolution:ils}{基因树}可能与\omterm{prop:b3:molecular-evolution:ils}{物种树}不同，
因此数据越多，错误答案反而越可信。溯祖法
对每一候选\omterm{prop:b3:molecular-evolution:ils}{物种树}计算期望的
\omterm{prop:b3:molecular-evolution:ils}{基因树}分布，并选出最能
解释所观测频率的\omterm{prop:b3:molecular-evolution:ils}{物种树}。
\textbf{24.} \omterm{prop:b3:molecular-evolution:ils}{不完全谱系分选}发生在所有现代人的共同祖先
种群中，因此它会使每个人类种群与尼安德特人的亲缘关系
相当。非洲以外人群中的过量
表明尼安德特人的等位基因是在非洲以外人群的祖先
离开非洲之后才进入的：即基因渗入，发生于约五万年前。
\textbf{25.} $\mu \approx 2.3\times 10^{-8}$ 每位点每代；
$\omega \approx 0.14$；约 \qty{30}{\%} 的\omterm{prop:b3:molecular-evolution:ils}{基因树}与
\omterm{prop:b3:molecular-evolution:ils}{物种树}不一致。
\end{solution}
